\cotefile{91}
\scanpage{2}

Représentabilité des Foncteurs \uuline{Hom}, $\prod_{T/S}X/T$ de Greenberg etc.
% archive sheet number: 1
\scanpage{3}

Représentabilité relative, \unsure{de|à} $\displaystyle\prod_{X/S}P/X$
(sans hyp. projectives)

\centerline{\rule{1.6cm}{0.3pt}}

``Flattening functor''

\centerline{\rule{1.6cm}{0.3pt}}

Passage au quotient

\centerline{\rule{1.6cm}{0.3pt}}

\[
\begin{aligned}
&\underline{\mathrm{Hom}}_{S}(X,Y)\\
&\underline{\mathrm{Aut}}_{S}(X)
\end{aligned}
\]
% archive sheet number: 2
\scanpage{5}

% Two parallel pen strokes stand in the upper left margin.
\uline{Th} \quad $f:X\longrightarrow S$ morphisme \inserted{\uline{plat}} de type fini, avec
$S$ loc noeth, $Y$ \deleted{\unsure{sous-préschéma} \unsure{fermé} de $X$}\inserted{$\longrightarrow X$ un \unsure{mono|morph.} de type fini}.

Pour que $\displaystyle\prod_{X/S}Y/X$ soit représentable, il f. et s. qu'il commute aux limites projectives adiques d'anneaux artiniens\inserted{, si $Y\to X$ est une \unsure{imm} (\unsure{imm} \uuline{f.})}, et alors il est représentable par un sous-préschéma (\uuline{fermé}) de $S$.
% slanted
\margin{il suffit que $Y\to X$ soit un \uuline{monomorphisme}}
\centerline{\rule{0.94\linewidth}{0.3pt}}
\deleted{\unsure{fpqc|qui}}

\uline{Dém}. Vérifions les conditions générales du th. de représentabilité par morphismes non ramifiés séparés

$1^\circ)$ $F$ faisceau pour fpqc : \deleted{\unsure{EGA IV 2|EGA IV 8}} clair EGA IV 2

$2^\circ)$ $F$ commute aux $\underrightarrow{\lim}$ d'anneaux : clair EGA IV 8
% slanted
\margin{$X$ t.f.}

$3^\circ)$ C'est la condition de commutation aux $\underleftarrow{\lim}$ adiques d'artiniens.

$4^\circ)$, $\deleted{\unsure{6}}5^\circ)$, $6^\circ)$ Dans le cas d'un monomorphisme $F\longrightarrow S$ se réduit à ceci

si $A\longrightarrow A'$ est un hom injectif d'artiniens \inserted{$A'/A$ fini}
alors $F(A)\to F(A')$ est bijectif, ou encore $F(A')\ne\emptyset\Longrightarrow F(A)\unsure{=|\ne}\emptyset$.
Il faut prouver que si $X_{A'}=Y_{A'}$, on a $X_A=Y_A$. C'est clair, car $X_A$ \deleted{\unsure{étant}} plat sur $A$,
$\mathcal O_{X_A}\deleted{\unsure{\otimes_A}}\to\mathcal O_{X_A}\otimes_A A'$ est injectif.
% slanted
\margin{$X$ \uline{plat}}

$7^\circ)$ Si $T$ sur $S$ est \deleted{\unsure{\illegible{1}|noeth.}} local, noeth, muni d'un idéal de carré nul $J$, $T_0=V(J)$, $U$ ouvert de $T$ tel que $T$ en soit l'adhérence schématique, \uline{alors} pour
$\deleted{\unsure{\xi\in F|\xi\in\mathcal F|F(T)}\,\unsure{tF}}$
% slanted
\margin{$X$ \uline{plat} \deleted{\unsure{plat}}}
\diagram{\begin{tikzcd}[ampersand replacement=\&]
F(T)\arrow[r,dash]\arrow[d] \& F(T_0)\arrow[d]\\
F(U)\arrow[r] \& F(U_0)
\end{tikzcd}}
est cartésien.
% archive sheet number: 4
\scanpage{6}

% The two lines are upside down in the scan.
\deleted{\uline{Théorème} \quad $f:X\longrightarrow S$ \inserted{plat} de présentation finie, $Y$ un \unsure{sous-préschéma|sous-présch.} fermé de $X$ de présentation finie. Pour}
% archive sheet number: 5
\scanpage{7}

\deleted{En d'autres termes}\inserted{\unsure{i.e.|r.e.|r.c.}}, si \inserted{lorsque $Y$ est \inserted{\unsure{déjà|supposé}} fermé dans $X$} on a $F(T_0)\ne\emptyset$, $F(U)\ne\emptyset$, alors $\deleted{\unsure{x|X}}$ $F(T)\ne\emptyset$. \inserted{En fait, il suffit ici $F(U)\ne\emptyset$ pour $F(T)\ne\unsure{0|\emptyset}$.}
\deleted{\unsure{D} si $X_{T_0}=Y_{T_0}$ et} \inserted{En d'autres termes si} $X_{\deleted{\unsure{T}}U}=Y_U$, alors $X=Y$.
En effet, comme $\deleted{\unsure{S}}T$ est l'adhérence schématique de $U$ et $X_T$ plat sur $T$, $X_{\deleted{\unsure{S}}T}$ est l'adhérence schématique de $X_U$, donc \deleted{on gagne} $Y_T=X_T$ puisque $Y_T$ sous-préschéma fermé de $X_T$ qui \unsure{majore|possède|majorant} $X_U$.

$8^\circ)$ Si $T$ sur $S$ irréductible muni de $J$ de carré nul, $T_0=V(J)$, si $F(T_0)\ne\emptyset$, \deleted{alors $\exists$ ouvert dense \unsure{tel|de} $F(\underline{\mathcal O}_{T,\unsure{y|\eta}})\ne\emptyset$} alors $\exists$ ouvert dense $U$ de $T$ \deleted{sur lequel $F$} tel que $F_U$ repr.

% slanted
\margin{$X$ plat t.f.}
En effet, il suffit que $\deleted{\unsure{F|X}}X_U$ soit ``ess. libre'' sur $U$, et on sait qu'on peut trouver $U$ de cette façon EGA IV 11

\deleted{\unsure{D}} De plus, on a vérifié $7^\circ$ sous la forme plus forte \unsure{signifiant} que le foncteur est \uuline{propre}.

\deleted{\uline{Corollaire} \unsure{So}}

\deleted{\uline{Corollaire} Supposons que $X$ soit un $S$-groupe, $Y$ un $S$ sous-groupe. Supposons que $X$ ait la propriété suivante :

Pour tout \deleted{\unsure{sch}} $T$ \inserted{sur $S$} spectr. d'un anneau de valuation discrète, \deleted{sur $S$,} et tout \deleted{sg \unsure{pr}} sous groupe \deleted{\unsure{fermé}} $\unsure{Z|H}$ de $X_T$ qui \deleted{\unsure{contienne|contient} la fibre fermée \unsure{entière} spéciale} \unsure{majore en|ou pour un}}
% archive sheet number: 6
\scanpage{8}

La commutation aux $\underleftarrow{\lim}$ adiques peut aussi s'exprimer ainsi ;

(L) Pour tt \deleted{\illegible{1}} $T$ \inserted{sur $S$} spectr. d'un anneau local noeth. \deleted{complet sur $S$}\inserted{(qu'on peut supposer \uline{complet}, à corps rés. alg. clos.)}, et tel que $Y_T$ majore un voisinage ouvert induit de la fibre spéciale de $X_T$, on a $\deleted{\unsure{Y_T|X_T}}X_T=Y_T$.

\uline{Corollaire} Cas $X$ propre sur $S$ ; \uuline{O.K.}

Si $X$ est un groupe, $Y$ un sous-groupe, on voit que $Y_T$ \deleted{contient}\inserted{majore} un voisinage ouvert induit de la section unité, donc $Y_T\to X_X$ est \inserted{un \uline{mono}} étale, \deleted{\unsure{donc}} étant une imm \deleted{fermée}, est une immersion ouvert\unsure{\deleted{e}|e} \deleted{et fermée}. \deleted{\unsure{Donc dans ce cas|Donc dans le cas}} \inserted{De plus, la vérification $Y_T=X_T$ \unsure{est la|est|se fait} } \deleted{même} chose ``ensembliste'', donc se ramène au cas où $T$ est le spectre d'un anneau de valuation discrète. Donc

\uline{Corollaire} Dans ce cas, la condition envisagée équivaut à la suivante : Pour tt $T$ sur $S$ spectr. d'un anneau \deleted{local noeth} \deleted{de \unsure{valuation}}\inserted{\uline{de valuation discrète}} (qu'on peut supposer complet à corps rés. alg clos), tel que $Y_T$ soit un sous-groupe ouvert et fermé de $X_T$ qui contient la fibre \deleted{générique}\inserted{spéciale}, on a $Y_T=X_T$. Cette condition est vérifiée dans chacun des cas suivants.

a) $X$ \uline{à fibres connexes}

b) \uline{Il existe une famille de ss-\deleted{groupes finis}\inserted{schémas}} $\deleted{\unsure{H}}G_i^{\unsure{?}}$ de $G$, finis sur $S$, et tels que pour chaque $s\in S$, la réunion ensembliste des $G_{i,\bar s}\cdot G_{\bar s}^{0}$ soit $G_{\bar s}$.
% archive sheet number: 7
\scanpage{9}

\deleted{\uline{Définition} Un \deleted{\unsure{groupe est dit|groupe}} schéma en groupes $X$ \inserted{de type fini} sur $S$ est dit \unsure{tto-ri} de}

\uline{Définition} D'une condition sur un préschéma en \unsure{groupe|groupes} $G$ de \deleted{\inserted{\unsure{prés|plat}} type \unsure{finie|fini}} sur $S$ : \deleted{\unsure{localement|localement noeth}} pour tt schéma local $T$ sur $S$, \deleted{\unsure{il|et}} tout sous-groupe ouvert \deleted{et fermé} $U$ de $G_T$ qui majore la fibre spéciale, on a $U=G_T$.
% slanted
\margin{\uline{Condition R}}
Lorsque \deleted{$G$ \unsure{sur $S$ est|est}} $G$ est de type fini sur $S$ loc. noeth, il suffit de tester aux $T$ spectr. d'un anneau de valuation discrète [complet, corps rés. alg. clos].

On trouve donc

\uline{Corollaire} Soit $\deleted{G}X$ de \deleted{type fini sur présentation}\inserted{type} finie sur $S$ \inserted{loc noeth}, $\deleted{H}Y$ un sous-groupe fermé de, \deleted{présentation finie}. Supposons que $X$ satisfasse R. Alors $\displaystyle\prod_{X/S}Y/X$ est représentable par un sous-préschéma \uline{fermé} de $S$.

\unsure{Exemples|\deleted{Exemples}}

\deleted{Soit}

\deleted{\uline{Corollaire} Soit $G$ schéma en groupes noeth sur $S$ loc. noeth, $H$ un sous-groupe fermé de type fini sur $S$, satisfaisant la condition R. Alors $\underline{\mathrm{Cent}}._{\unsure{G|B}}(H)$ \inserted{(si $G$ séparé sur $S$)} et $\underline{\mathrm{Norm}}_G(H)$ sont représentables par un ss-préschéma fermé de $G$. De $=$ $\mathrm{Transp}_G(H,K)$}
% archive sheet number: 8
\scanpage{10}

\uline{Corollaire} Soit $u:H\to G$ un hom de présch en groupes \inserted{sur $S$}, avec $G$ loc. noeth. \inserted{et \deleted{séparé sur $S$}}, $H$ de type fini sur $S$, \deleted{Alors} $H$ \inserted{plat et} satisfaisant la condition R (p. ex. à fibres connexes...). \deleted{\inserted{Si $G$ séparé sur $S$}} Alors $\underline{\mathrm{Cent}}_G(u)$ est représentable par un sous-préschéma \deleted{fermé} de $G$, \deleted{\unsure{Si de|Celui-ci}} fermé si $G$ est séparé sur $S$.

Si $K$ est un sous-groupe de $G$, \deleted{alors}\deleted{\unsure{de t.f.}} \inserted{de t.f. sur \unsure{$K$|$S$}} $\underline{\mathrm{Transp}}_G(\unsure{u|u'},K)$ est représentable, et si $K$ est un sous-préschéma de $G$ (resp. sous-préschéma fermé de $G$) il en est de $=$ de $\underline{\mathrm{Transp}}_G(u,K)$. Enfin, si \deleted{\unsure{en tt|$u$ et $H$}} $K$ est également plat \deleted{\unsure{et}} sur $S$ et satisfait la condition R, \deleted{alors $\mathrm{Transp}_G$} et $u$ un \uline{mono}, alors $\deleted{\mathrm{Transp}_G(u,K)}$ $\mathrm{Transp}.\deleted{\unsure{\mathrm{str}}}_G(H,K)$ est représentable par un \deleted{\unsure{sous-pr|sous-pré}} schéma de type fini sur $G$, et c'est un sous-préschéma (resp sous-préschéma fermé) s'il en est de $=$ de $H$ et $K$ dans $G$.

En particulier, $\underline{\mathrm{Norm}}_G(H)$ est représentable, c'est un sous-préschéma (sous-préschéma fermé) \inserted{de $G$,} si $H$ l'est. \deleted{\illegible{2}}
% archive sheet number: 9
\scanpage{11}

\uline{Variantes} \unsure{débarrasser|débarrassées} d'hypothèses noethériennes...

\uline{Variante} Pour un morphisme $Y\to X$ \uline{non ramifié}.

\uline{N. B.} Il faudrait également \uline{moduler} le th. précédent ....

\centerline{\rule{2.6cm}{0.3pt}}
\centerline{\rule{2cm}{0.3pt}}
\centerline{\rule{0.8cm}{0.3pt}}

\uline{Théorème} Soit $f:X\unsure{\longrightarrow|\text{---}}S$ morphisme de type fini, $S$ loc noeth, $E$ cohérent sur $X$. Soit $F$ sur $S$ le ss-foncteur de $S$
\[
F(S')=\begin{cases}
\emptyset & \text{si }F_{S'}\ \text{non plat }\inserted{\text{(resp. non fid plat)}}\text{ sur }S'\\
\{\emptyset\}&\text{sinon}
\end{cases}
\]
Pour que ce foncteur soit représentable, il f. et s. \unsure{qu il|qu'il} commute aux limites projectives adiques d'anneaux artiniens.

Vérification des divers \unsure{pts|pts.}

$1^\circ)$ Faisceau pour fpqc Clair

$2^\circ)$ Commute aux $\underrightarrow{\lim}$ d'anneaux : EGA IV 8 et \deleted{\unsure{IV}} 11
% slanted
\margin{$X$ type fini sur $S$, $F$ \unsure{noeth|noeth.}}

$3^\circ)$ Condition de commutation aux $\underleftarrow{\lim}$ d'artiniens

$4^\circ)$ $5^\circ)$ $6^\circ)$ Si $A\to A'$ injectif ($A'$ fini sur $A$ local artinien) alors $F(\deleted{\unsure{A'}}A)\to F(A')$ bijectif i.e $F(A')\ne\emptyset\Rightarrow F(A)\ne\emptyset$ i.e. $X_{A'}$ plat (fid plat) sur $A'$ $\Longrightarrow X_A$ plat (fid plat) sur $A$ \quad EGA IV 11

$7^\circ)$ \unsure{on a|\ldots} $F(T_0)\ne\emptyset$, $F(\deleted{\unsure{H}}U)\ne\emptyset\Longrightarrow F(T)\ne\emptyset$ i.e. $X_{T_0}$ plat et $X_U$ plat $\Longrightarrow X_T$ plat \quad EGA IV 11.
% archive sheet number: 10
\scanpage{12}

$8^\circ)$ .... Trouver $U$ ; il suffit encore que $X_{U_0}$ soit ``ess. libre'' sur $U_0$, OK par EGA IV 11

La condition énoncée dans le th. signifie ceci :

(L) Pour tt $T$ spectr d'un anneau de \deleted{valuation discrète}\inserted{local noeth}, (complet à corps résiduel alg. clos si on y tient, \deleted{\unsure{tout|tel}} tel que $F_T$ sur $X_T/T$ soit plat sur $T$ dans un voisinage de la fibre générique [resp. et $F_{T_0}\ne0$, où $T_0$ le pt fermé...], \deleted{alors} $F_T$ soit plat sur $T$, \deleted{\unsure{Lorsque|Exemple}}

% slanted
\margin{\uline{N.B.} On peut supposer $\dim T=1$ \deleted{$T$ irréductible} $T$ \unsure{irréduc et|intègre} i.e. $\mathrm{Ass}\,\mathcal O_T$ réduit à un élément, et $T_{\mathrm{réd}}$ spectre d'un anneau de valuation discrète...}

\uline{Corollaire 1} O.K. si $f:X\to S$ est \uline{propre}.

\uline{Corollaire 2} O.K. si \deleted{\unsure{G est|G}} $F=\mathcal O_X$, $X$ un préschéma \uline{en gpes} sur $S$, satisfaisant la condition R plus haut ....

\uline{Corollaire 3} Soit $f:X\to S$ un morphisme de \deleted{t fini, $S$ loc noeth.}\inserted{prés. finie} Conditions équivalentes

(i) $f$ se factorise en $f'f''$ : \deleted{\inserted{$X\to S'\to S$, $S'$ de t.f. sur $S$,}} avec $f''$ fid plat, $f'$ \uuline{\unsure{mm|imm|mono|imm.}}

(i bis) \rule{4cm}{0.3pt} avec $f''$ plat, $f'$ mm,

(ii) \unsure{La|Le} \deleted{\unsure{sa}} $X\times_S X\xrightarrow{\mathrm{pr}_1}X$ est \uline{plat}
% Two vertical pen strokes beside (i)–(ii).
% archive sheet number: 11
\scanpage{13}

Sous ces conditions, \unsure{la|\deleted{la}} factorisation (i) est unique à isom près, et $S'$ est de \deleted{t.f.}\inserted{prés. finie} sur $S$.

(i) $\Rightarrow$ (i bis) $\Rightarrow$ (ii) trivial. De $=$ l'unicité de la factorisation \inserted{(i)}, car $X\to S'$ sera \unsure{qu. cpt|quotient}, donc couvrant pour fpqc donc $S'=X/R$, où $R=X\times_S X$.
Reste à prouver (ii) $\Rightarrow$ \inserted{le faisceau} $X/R$ est représentable par un $S'$ de \inserted{\unsure{loc.|prés.}} type fini sur $S$. \deleted{(N.B. alors, cela quitte à \unsure{part} sur $S$ car $X\to S'$ est}

[N.B. Alors $X\to S'$ est couvrant pour \unsure{fppc|fpqc|fppf}, comme après ch. \unsure{mt} de base $X\to S'$ il donne $X\times_S X=R\longrightarrow X$ qui est fid plat \unsure{par hypoth|qu. cpt}, il est fid plat qu. cpt, lui-$=$, ce qui implique \inserted{aussi} que $S'$ est \inserted{qu. cpt donc} de t.f. \deleted{\unsure{donc qu.|de présentation}} sur $S$].

Or \deleted{\unsure{comme|le foncteur} $X/R\to S$} le foncteur représenté par $X/R\hookrightarrow S$ est \unsure{ici|i.e.} le suivant :
\[
F(\deleted{\unsure{S}}T)=\begin{cases}
\emptyset&\text{si } X_T\ \text{non fid plat sur }T\ \inserted{[\Longleftrightarrow X_T\to T\ \text{couvre fppc}]}\\
\{\emptyset\}&\text{sinon}
\end{cases}
\]
D'ailleurs on est ramené au cas $S$ affine, \uuline{noethérien}. \deleted{et} Alors on applique le Th, qui nous dit qu'il suffit de vérifier \deleted{que} (L). Or sous la condition de cette condition, $X_T\to T$ est couvrant [car il existe un pt $x$ de la fibre spéciale, et $\unsure{\mathrm{Sp}|\mathrm{Spec}}\,\mathcal O_{X,x}$ \unsure{--|$\to$} $\mathrm{Sp}\,\mathcal O_{T,t_0}$ est fid plat] donc on gagne !
% archive sheet number: 12
\scanpage{14}

\uline{Exemple}

\uline{Corollaire} Soit $G$ un préschéma en groupes de \deleted{type}\inserted{prés} \unsure{fini|finie} sur $S$, \deleted{loc. noeth.} $P$ un \deleted{fibré principal} préschéma \unsure{de|\deleted{de}} prés. finie sur $S$ dans lequel $G$ opère, de façon que $P$ soit form. pr. hom. sous $G$. \unsure{Dans|Sous} ces conditions, il existe un mono de préschémas $S'\hookrightarrow S$, tel que $P\longrightarrow S$ se factorise par $P\unsure{\text{---}|\to}S'$, et que \deleted{\unsure{Reste} $P'$} $P$ devienne \deleted{formellement} principal hom sous $G_{S'}$ i.e $P\to S'$ est \deleted{plat} \uline{fid.} \uline{plat}. [De \deleted{\unsure{tt}} plus, $S'$ est \unsure{uniqt|unique,|uniquement|unique} déterminé à $S$-isom près, et est de prés. finie sur $S$].

\uline{Corollaire} Soit $u:G\unsure{\to|\text{---}}H$ un hom de $S$-préschémas en groupes, \deleted{de présentation} $u$ de présentation finie. Supposons $\mathrm{Ker}\,u=\deleted{\unsure{K}}N$ [qui est de prés finie sur $S$] \uuline{plat} sur $S$. \uline{Alors} \deleted{\unsure{$u$ se factorise|le foncteur}} $G/N$ \uline{existe}.
% A vertical pen line marks the second corollary.
% archive sheet number: 13
\scanpage{16}

\diagram{\begin{tikzcd}[ampersand replacement=\&]
 Z \arrow[r,"\varphi"] \arrow[dr,swap,"{\deleted{\unsure{\varphi|f}}h}"] \& X \arrow[d,dash,swap,"f"] \& Y \arrow[dl,dash,"g"] \& Z' \arrow[r,"\varphi'"] \arrow[rr,bend left=17,"{v}"] \arrow[dr,swap,"h'"] \& X' \arrow[r,swap,"u"] \arrow[d,swap,"f'"] \& Y' \arrow[dl,"g'"]\\
 \& S \& \& \& S' \&
\end{tikzcd}}

\margin{Énoncer plutôt en termes de $\prod_{X/S}P/X\longrightarrow\prod_{Z/S}Q/Z$}

$f,h$ plat et propres sur $S$ loc. noeth

$g$ \deleted{$\unsure{g,\varphi|g,h}$} de type fini [et séparé?]
\margin{inutile si $Z\to X$ est \uline{surjectif}}
% slanted

\margin{\uline{NB.} l'hypothèse $S$ loc. noeth peut être \unsure{remplacé|remplacée} par des hypothèses de présentation finie sur $f,g,h$.}

Pour $\varphi$ on a
\[
\left|\begin{aligned}
&u \qquad \unsure{\mapsto u\circ\varphi|\text{res}\,u\circ\varphi}\\
&\underline{\mathrm{Hom}}_S(X,Y)\longrightarrow\underline{\mathrm{Hom}}_S(Z,Y)
\end{aligned}\right.
\]

a) Le sous-foncteur de $\underline{\mathrm{Hom}}_S(X,Y)$ où $\underline{\mathrm{Hom}}_S(X,Y)$ est \uline{non ramifié} sur \deleted{$\unsure{S}$}\inserted{$\underline{\mathrm{Hom}}(Z,Y)$}, est un sous-foncteur ouvert.

Il correspond à $u':X'\to Y'$ tel que, pour tout $s\in S'$, $\mathrm{Hom}(\unsure{u_s^{**}|u_s^{\prime *}}(\Omega^1_{\deleted{\unsure{X}}Y'_s/k(s)}),k(s))\unsure{-|\to}\mathrm{Hom}(\unsure{w_s^{**}|w_s^{\prime *}|v_s^{*}}(\Omega^1_{Y'_s/k(s)}),k(s))$ soit injectif. \uline{Or} on peut trouver sur $S'$ des Modules \uline{cohérents} $M_u,M_w$ tels que pour \deleted{\unsure{tt|un|tout}} $J$ qu. coh. sur $S'$, on ait des isom fonctoriels
\[
\begin{aligned}
\mathrm{Hom}_{\mathcal O_{X'}}(u^*(\Omega^1_{Y'/S'}),\unsure{f|\mathcal O}_{X'}\mathbin{\unsure{\otimes|{}}}_{\mathcal O_{S'}}J)&\simeq\mathrm{Hom}_{\mathcal O_S}(M_u,J)\\
\mathrm{Hom}_{\mathcal O_{Z'}}(w^*(\Omega^1_{Y'/S'}),\mathcal O_{Z'}\otimes_{\mathcal O_{S'}}J)&\simeq\mathrm{Hom}_{\mathcal O_S}(M_w,J)
\end{aligned}
\]
\margin{NB \uline{réduction} \uline{immédiate} \uline{au cas} $S'$ \uline{noethérien}}
% slanted
et la formation de $M_u,M_w$ est compatible avec l'extension de la base. Ceci dit, on a un hom naturel
\[
M_w\longrightarrow M_u\tag*{(*)}
\]
qui définit par $\mathrm{Hom}(\ ,J)$ l'hom naturel
\[
\mathrm{Hom}_{\mathcal O_{X'}}(u^*(\Omega^1_{Y'/S'}),\mathcal O_{X'}\otimes_{\mathcal O_{S'}}J)\longrightarrow
\mathrm{Hom}_{\mathcal O_{Z'}}(w^*(\Omega^1_{Y'/S'}),\mathcal O_{Z'}\otimes_{\mathcal O_{S'}}J)
\]

Dire \deleted{\unsure{alors}}\inserted{que} la condition de non ramification est satisfaite en un pt $s\in S'$, signifie qu'en ce point
\[
M_w\otimes k(s)\longrightarrow M_u\otimes k(s)
\]
est \uline{surjectif}, et par Nakayama ceci restera vrai en
% archive sheet number: 15
\scanpage{17}

voisinage.

Nous désignerons par
\[
F=F_\varphi=\underline{\mathrm{Hom}}_S(X,Y;\varphi)
\]
\unsure{le|les} sous-foncteur précédent de $\underline{\mathrm{Hom}}_S(X,Y)$, dont la valeur en $S'$ est l'ens des $S'$-hom $u:X'\to Y'$ tels que, pour $v=u\varphi'$, et considérant $M_{\unsure{v|w}}^{\unsure{\cdot|\prime|{}}},M_u,M_v\to M_u$ correspondants, $M_v\to M_u$ soit surjectif.

\uline{Théorème} $\overset{\frown}{\underline{\mathrm{Hom}}_S(X,Y;\varphi)}\longrightarrow\overset{\frown}{\underline{\mathrm{Hom}}_S(Z,Y)}$ est représentable par morphismes non ramifiés.

\uline{Corollaire} Si $\underline{\mathrm{Hom}}_S(Z,Y)$ est représentable, alors $\underline{\mathrm{Hom}}_S(X,Y;\varphi)$ l'est également, et le morphisme $\underline{\mathrm{Hom}}_S(X,Y;\varphi)\longrightarrow\underline{\mathrm{Hom}}_S(Z,Y)$ \unsure{est|\deleted{est}} non ramifié.
% Vertical marginal rule alongside the last two lines of the corollary.

\uline{Principe de la démonstration} On peut supposer donné un $v:\mathrm{Hom}_S(Z,Y)$, et on doit représenter le sous-foncteur de $\mathrm{Hom}_S(X,Y)$ dont la valeur en $S'$ est l'ens des $S'$-morphismes $\unsure{u'|u^{\bullet}}:X'\to Y'$ tels que 1°) $u'\circ\varphi'=v'$ \quad 2°) la condition $M_{v'}\to M_{u'}$ surjectif. On applique le critère principal de représentabilité \inserted{pour} des \deleted{foncteurs}
\margin{N.B. Réduction au cas $S$ noeth. \unsure{immédiate|immédiats|irréductible}, \unsure{dans suite|dans doute|dans toute question} sur foncteurs rel. représentables}
% slanted
% archive sheet number: 16
\scanpage{18}

préschémas \deleted{loc de type fini} \inserted{non ramifiés} séparés sur $S$. Seules les \inserted{deux dernières} conditions de modération se vérifient par des techniques habituelles. \deleted{\unsure{On a|On ne}} \unsure{Celles|Celle} au point générique ne doit pas \unsure{faire|\deleted{faire}} d'ennui, avec le calcul habituel d'obstructions.
% An angular ink line separates the first assertion from the line about the generic point; a wavy stroke runs under the latter.
\margin{\uline{Tout} \uline{\unsure{serait à|suit en|se lit en}} \uline{\unsure{simplifier|simplifiant}}}
% slanted

Reste le cas où $S$ est local noeth irréductible de dim 1, avec Idéal $J$ annulé par le nilradical, d'où $S_0=\mathrm V(J)$, et où on dispose de $u_0:X_0\to Y_0$ (tel que $u_0\varphi_0=v_0$) et un prolongement $u_{\unsure{y|\eta}}:\deleted{\unsure{X}}X_y\to Y_y$ ($y$ le pt générique de $S$) tel que $u_y\varphi_y=v_y$, $u_{y0}=u_{0y}$. Question de prolonger $u_0$ en $u$ tel que $u\varphi=v$, \deleted{\unsure{u}} Soient $P(u_0,S),P(u_0,S_y);P(v_0,S),P(v_0,S_y)$ les prolongements de $u_0$ resp. $v_0$ \deleted{sur} $S$, \deleted{\unsure{resp|et}}\inserted{et} de $u_{0y}$ resp. $v_{0y}$ sur $S_y$, d'où le diagramme commu
% Vertical marginal rule beside the question about prolongation.
\diagram{$(1)$\quad\begin{tikzcd}[ampersand replacement=\&]
 P(u_0,S)\arrow[r]\arrow[d]\&P(u_0,S_y)\ {*}\arrow[d]\\
 {*} \ P(\deleted{\unsure{u}}v_0,S)\arrow[r,dash]\& P(v_0,S_y)
\end{tikzcd}}

d'ens \unsure{form.|form} princ. homogènes, sous les groupes
\diagram{$(1')$\quad\begin{tikzcd}[ampersand replacement=\&]
 \mathrm{Hom}(M^{\deleted{\unsure{v|*}}}_{u_0},J)\arrow[r]\arrow[d]\&\mathrm{Hom}(M_{u_0},J_y)\arrow[d]\\
 \mathrm{Hom}(M_{v_0},J)\arrow[r]\&\mathrm{Hom}(M_{v_0},J_y)
\end{tikzcd}}

Par hypothèse on tient des éléments aux coins
% archive sheet number: 17
